\ifdefined\isMainDocument
\else
    \documentclass{article}
    \usepackage{fullpage}
    % \usepackage{geometry}
    % \geometry{top=0.5in, bottom=0.5in, left=0.5in, right=0.5in}
    
    \usepackage{wrapfig}
    \usepackage{lineno}
    \usepackage{amsmath, amsthm, amsfonts, amssymb, amscd}
    \usepackage{bm, bbm, mathrsfs}
    \usepackage{mathtools}
    \usepackage{nicefrac}
    \usepackage{caption}
    %\usepackage{natbib}
    \usepackage[style=numeric, sorting=none]{biblatex}
    \addbibresource{refs.bib}
    \usepackage[colorinlistoftodos]{todonotes}
    \usepackage{graphicx}
    \graphicspath{figures/}
    \usepackage{setspace}
    \setstretch{1.5}

    \begin{document}
    \linenumbers
\fi

\section{Materials and Methods}
\subsection{Empirical Pedigree Parameters and Variance Posteriors}
Estimates of additive genetic variance in relative fitness ($V_A$) were compiled from published pedigree-based animal model analyses of longitudinally monitored wild vertebrate populations. Studies were included if they: (i) reported $V_A$ for a composite fitness measure or a major fitness component (annual or lifetime reproductive success, survival to reproductive age, or a multivariate index thereof); (ii) provided sufficient pedigree depth and sample size to support reliable separation of additive from residual environmental variance; and (iii) reported posterior credible intervals or standard errors for the $V_A$ estimate. All estimates were expressed as additive genetic variance in fitness on the relative fitness scale ($\bar{w} = 1$), consistent with the parameterization of Equations (1) and (5); published estimates reported on a different scale were standardized prior to analysis [author: describe standardization procedure]. Where multiple published estimates were available for the same population, [author: state resolution rule—most recent, most comprehensive pedigree, or population-level average]. The full dataset, with source studies and population identifiers, is provided in [Table/Supplementary Table X].

\subsection{Mating System Parameterization}
For each species, the extra-pair paternity rate $\alpha$ was obtained from published molecular parentage analyses of the same or a conspecific monitored population [author: specify source for each species and confirm whether $\alpha$ is defined as the proportion of offspring or the proportion of broods containing at least one extra-pair offspring; the theoretical definition in Equations (3) and (4) is the proportion of offspring, and any conversion should be noted here]. For species lacking a direct estimate of $\alpha$, [author: state whether strict monogamy was assumed or whether those species were excluded from the mating-system-corrected prediction].
The mating-system scalar $\kappa(\alpha)$ was computed under both the random extra-pair siring model (Equation 3) and the fitness-proportional extra-pair siring model (Equation 4), and predictions are reported under both to bracket the effect of the assumed extra-pair mechanism.

\subsection{Genomic Architecture and Functional Target Size}
For species with chromosome-level genome assemblies, reference sequences and gene annotation files (GTF or GFF format) were retrieved from NCBI. Analysis was restricted to fully assembled autosomes with unambiguous correspondence to a major linkage group in the published linkage map (see below); unplaced scaffolds, unlocalized contigs, sex chromosomes, and mitochondrial sequence were excluded. Each assembled autosome was assigned to its corresponding linkage group on the basis of marker content and ordering reported in the linkage map source publication.

The between-chromosome partition of additive fitness variance was computed from annotated protein-coding gene counts following Equation (19):
\begin{equation*}
    \hat{f}i = \frac{G_i}{G_{\text{total}}}
\end{equation*}
where $G_i$ is the number of annotated protein-coding loci on autosome $i$ and $G_{\text{total}} = \sum_i G_i$ over all included autosomes. Gene counts were taken from the feature tables of NCBI assembly annotation reports. For species without available genome assemblies, $f_i$ was approximated as proportional to physical chromosome length, $f_i \approx L_i / L_{\text{total}}$.

For the mutational influx calculation (below), total genome-wide coding sequence length $L_{\text{coding}}$ was estimated independently of the gene-count partition. All CDS features in the annotation file were extracted and overlapping intervals on each strand collapsed into non-redundant contiguous ranges to prevent multiple isoforms and overlapping reading frames from inflating the functional target; the genome-wide collapsed coding length is $L_{\text{coding}} = \sum_i L_{\text{coding},i}$. This quantity enters only the $U_d$ calculation and does not affect the $\hat{f}_i$ partition.

\subsection{Genetic Map Lengths}
Genetic map lengths ($M_i$, in Morgans) were taken from published empirical linkage maps for each species. Where sex-specific and sex-averaged map lengths were both available, the sex-averaged value was used to reflect the expected map length across both parental gametes. Autosomes without a resolved major linkage group in the published map were excluded from the linked-constraint calculation; the selective variance apportioned to excluded chromosomes by $\hat{f}_i$ was redistributed to the unlinked background [author: verify this is the procedure applied, or describe the alternative].

For species lacking per-chromosome map data, $M_i$ was approximated as $M_i = \bar{\rho} \times L_i \times 10^{-8}$, where $\bar{\rho}$ is the published genome-wide average recombination rate in cM per Mb and $L_i$ is the physical chromosome length in base pairs.

The closed-form expressions for $Q^2(y)$ and $\overline{Q^2}_i$ (Equations 8 and 9) were derived under the linear approximation $r \approx |x - y|$, which is accurate for $|x - y| \ll 0.5$ Morgan. For the vertebrate chromosomes analyzed here, the dominant contribution to $\overline{Q^2}_i$ arises from closely linked locus pairs where this approximation is most accurate. For chromosomes substantially longer than 0.5 Morgan, the linear approximation overestimates the recombination fraction between distant loci relative to the Haldane mapping function, causing a mild underestimate of $Q^2$ from those pairs; because this produces a conservative (upward) bias in the predicted $N_e/N$, the approximation errs in the direction of understating constraint. The magnitude of this bias across the chromosome-length distribution of each species is quantified in the Supplementary Material.

\subsection{Mutational Parameters and the Mutation-Selection Balance Test}
The per-generation survival rate of selective variance, $Z = 1 - V_m / V_A$, requires an estimate of the mutational input $V_m \approx U_d s^2$ (Equations [X] and [X] in Theory). The genome-wide diploid deleterious mutation rate was estimated as:
\begin{equation*}
    U_d = 4\mu L_{\text{coding}}
\end{equation*}
where $\mu$ is the published per-base-pair per-generation point mutation rate for each species and the factor of four converts from a per-haplotype per-base-pair rate to a diploid rate over both coding and regulatory sequence. The regulatory target is approximated as equal in length to the annotated exome, following empirical estimates of the ratio of constrained non-coding to coding sequence in vertebrate genomes [cite: e.g., Meader et al. 2010 or equivalent]; this multiplier doubles the effective functional target relative to coding sequence alone and enters the calculation as the factor of two in $U_d = 2\mu \times 2L_{\text{coding}}$.
Because the distribution of fitness effects of deleterious mutations is not independently constrained for the non-model vertebrate populations studied here, we do not fix a point estimate of $s$. Instead, the implied average heterozygous selection coefficient $s_d$ was inferred by solving $V_m = U_d s_d^2$ for $s_d$ given the empirically estimated $V_m = V_A(1 - Z)$ at each value of $V_A$. This was evaluated across the joint 95\% credible intervals of the published $V_A$ posterior and the empirical uncertainty in $\mu$, enabling a test of whether the maintenance of observed standing variance requires selection coefficients consistent with known vertebrate DFE parameters and tolerable levels of mutational load.

\subsection{Empirical Nucleotide Diversity}
\textit{[Author input required. This subsection must describe: (i) the source of genome-wide or synonymous-site nucleotide diversity $\pi$ for each species in the dataset, including the specific studies, variant call sets, or databases consulted; (ii) the genomic compartment used for the diversity estimate (e.g., fourfold-degenerate sites, synonymous sites, intergenic sequence, or genome-wide); (iii) the operational definition of $N_e/N$ used for empirical comparison—whether computed as $\pi/(4\mu)$ divided by census size $N$, or derived from a directly reported $N_e$ estimate; and (iv) how census size $N$ was estimated or obtained for each monitored population, since $N_e/N$ is the predicted quantity and $N$ enters the denominator. Without this subsection the theoretical predictions in Equations (11) and (13) cannot be connected to any empirical observation, which is the central comparison of the paper.]}

\ifdefined\isMainDocument
\else
    \printbibliography
    \end{document}
\fi